% ==================================================================
%  WEEK 5
% ==================================================================
\chapter{Service Level Agreements, Cloud Economics, and Green Computing}

\section{Introduction and Week Overview}
Week 5 investigates the operational, contractual, and environmental dimensions of cloud infrastructure. The core pillars covered in this module include:
\begin{enumerate}[leftmargin=*]
  \item \textbf{Service Level Agreements (SLAs) and Availability Modeling:} Contractual metrics, Service Level Objectives (SLOs), uptime and downtime mathematics, and tiered service credit penalty formulations.
  \item \textbf{Cloud Economics and Utility Pricing:} Total Cost of Ownership (TCO), peak-to-average demand analysis, utility premium thresholds ($U < P/A$), and penalty integrals under dynamic workloads.
  \item \textbf{Big Data Processing with Hadoop and HDFS:} Distributed file storage architectures, input split calculations, NameNode metadata management, and MapReduce execution mechanics (Partitioner vs Combiner).
  \item \textbf{Green Cloud Computing and Energy Efficiency:} Data center power consumption breakdown, dynamic server energy behavior, Power Usage Effectiveness (PUE), Data Center Infrastructure Efficiency (DCiE), and Dynamic Voltage and Frequency Scaling (DVFS).
\end{enumerate}

\section{Service Level Agreements (SLAs) in Cloud Computing}
\subsection{Contractual Foundations and Definitions}
\begin{defbox}{Service Level Agreement (SLA)}
A \textbf{Service Level Agreement (SLA)} is a legally binding formal contract negotiated between a Cloud Service Provider (CSP) and a Cloud Consumer. It codifies the minimum expected quality of service, availability guarantees, performance baselines, measurement methodologies, and financial liabilities (service credits/penalties) incurred when guarantees are breached.
\end{defbox}

An SLA is composed of discrete, quantifiable criteria known as \textbf{Service Level Objectives (SLOs)}:
\begin{equation}
\text{SLO}_i = \langle \text{Service Metric}, \text{Comparison Operator}, \text{Target Value}, \text{Measurement Window} \rangle
\end{equation}
Typical SLO parameters include:
\begin{itemize}[leftmargin=*]
  \item \textbf{Availability (Uptime):} Percentage of time the service responds to valid HTTP requests with non-error status codes.
  \item \textbf{Latency / Response Time:} Round-trip processing delay for API calls (e.g., $99^{\text{th}}$ percentile latency $< 150\text{ ms}$).
  \item \textbf{Throughput:} Number of discrete transactions or input/output operations per second (IOPS).
  \item \textbf{Data Durability:} Probability of preventing permanent data loss over a one-year period (e.g., AWS S3 guarantees $99.999999999\%$ ``eleven nines'' durability).
\end{itemize}

\subsection{Availability Mathematics and ``Nines'' Notation}
Service availability $A$ is mathematically defined as the ratio of scheduled operational uptime to total scheduled duration:
\begin{equation}
A = \frac{\text{MTBF}}{\text{MTBF} + \text{MTTR}} = \frac{T_{\text{scheduled}} - T_{\text{downtime}}}{T_{\text{scheduled}}} \times 100\%
\end{equation}
where $\text{MTBF}$ is the \textbf{Mean Time Between Failures} and $\text{MTTR}$ is the \textbf{Mean Time To Repair}.

\begin{center}
\begin{tabularx}{\textwidth}{>{\bfseries}l l X X}
\toprule
Availability Level & Percentage & Allowable Downtime / Year & Allowable Downtime / Month (30 Days)\\
\midrule
Two Nines & $99.0\%$ & $87.6$ hours ($3.65$ days) & $7.20$ hours\\
Three Nines & $99.9\%$ & $8.76$ hours & $43.8$ minutes\\
Four Nines & $99.99\%$ & $52.56$ minutes & $4.38$ minutes\\
Five Nines & $99.999\%$ & $5.26$ minutes & $26.3$ seconds\\
\bottomrule
\end{tabularx}
\end{center}

\subsection{Scheduled Operational Hours Downtime Formulation}
In enterprise contexts, customer applications frequently run on partial daily schedules rather than continuously. When an SLA availability guarantee $A_{\text{target}}$ applies to a service operating $h$ hours per day across $D$ calendar days:
\begin{equation}
T_{\text{scheduled}} = h \times D
\end{equation}
\begin{equation}
T_{\text{allowable downtime}} = (1 - A_{\text{target}}) \times T_{\text{scheduled}}
\end{equation}

\begin{solved}{SLA Outage Evaluation (Lecture Tutorial Problem 1)}
\Given A cloud provider guarantees $99\%$ service availability. A third-party client application runs for $h = 12\text{ hours/day}$ over a $D = 30\text{-day}$ billing month. An audit reveals a total outage of $T_{\text{actual}} = 10.75\text{ hours}$. Determine if the provider violated the SLA.
\Formula
\begin{equation}
T_{\text{allowable}} = (1 - 0.99) \times (12 \times 30) = 0.01 \times 360 = 3.6\text{ hours}
\end{equation}
\Steps
\begin{enumerate}[leftmargin=*]
  \item Compute total scheduled operational time: $T_{\text{scheduled}} = 12 \times 30 = 360\text{ hours}$.
  \item Compute allowable outage duration: $0.01 \times 360 = 3.6\text{ hours}$.
  \item Compare actual outage against threshold: $T_{\text{actual}} = 10.75\text{ hours} > 3.6\text{ hours}$.
  \item Measured availability achieved:
  \begin{equation}
  A_{\text{actual}} = \frac{360 - 10.75}{360} \times 100\% = \frac{349.25}{360} \times 100\% \approx 97.01\%
  \end{equation}
\end{enumerate}
\Answer The SLA guarantee was severely violated ($97.01\% < 99.0\%$). Actual downtime exceeded allowable downtime by $7.15\text{ hours}$.
\end{solved}

\subsection{Service Credit Penalty Models}
CSPs do not pay cash damages for SLA breaches; instead, they issue \textbf{Service Credits} that discount future billing cycles according to a tiered schedule:

\begin{solved}{Comprehensive SLA Credit Settlement (Lecture Tutorial Problem 2)}
\Given A company contracts a cloud service from Provider P for a $D = 30\text{-day}$ period with:
\begin{itemize}
  \item Target availability: $99.95\%$.
  \item Maximum operational hours: $h = 12\text{ hours/day}$.
  \item Base daily billing rate: $\$50/\text{day}$ (Total nominal fee: $30 \times \$50 = \$1500$).
  \item Tiered Service Credit policy:
    \begin{itemize}
      \item $\text{Uptime} < 99.95\% \implies 10\%$ credit discount.
      \item $\text{Uptime} < 99.00\% \implies 25\%$ credit discount.
    \end{itemize}
  \item Recorded outages across 5 separate incidents:
    $5\text{ h}$, $30\text{ min}$ ($0.5\text{ h}$), $1\text{ h } 30\text{ min}$ ($1.5\text{ h}$), $15\text{ min}$ ($0.25\text{ h}$), and $2\text{ h } 25\text{ min}$ ($2.4167\text{ h}$).
\end{itemize}
\Formula
\begin{equation}
T_{\text{outage}} = \sum_{k=1}^{5} t_k, \quad A = \frac{T_{\text{scheduled}} - T_{\text{outage}}}{T_{\text{scheduled}}} \times 100\%
\end{equation}
\Steps
\begin{enumerate}[leftmargin=*]
  \item Total scheduled service duration:
  \begin{equation}
  T_{\text{scheduled}} = 12 \times 30 = 360\text{ hours}
  \end{equation}
  \item Sum of all discrete outage durations:
  \begin{equation}
  T_{\text{outage}} = 5.0 + 0.5 + 1.5 + 0.25 + 2.4167 = 9.6667\text{ hours} \quad (9\text{ h } 40\text{ min})
  \end{equation}
  \item Calculate realized monthly uptime percentage:
  \begin{equation}
  A = \frac{360 - 9.6667}{360} \times 100\% = \frac{350.3333}{360} \times 100\% \approx 97.315\%
  \end{equation}
  \item Evaluate penalty tier:
  Since $97.315\% < 99.00\%$, the violation triggers the higher \textbf{25\% Service Credit}.
  \item Compute effective billing settlement:
  \begin{equation}
  \text{Nominal Cost} = 30 \times \$50 = \$1500
  \end{equation}
  \begin{equation}
  \text{Service Credit Awarded} = 0.25 \times \$1500 = \$375
  \end{equation}
  \begin{equation}
  \text{Effective Payable Cost} = \$1500 - \$375 = \$1125
  \end{equation}
\end{enumerate}
\Answer Effective amount payable is \textbf{\$1,125.00} (with \$375 awarded as SLA penalty credits).
\end{solved}

\section{Cloud Economics and Utility Pricing}
\subsection{Economic Drivers of Cloud Adoption}
Traditional IT infrastructure demands that enterprises size on-premises data centers to satisfy \textbf{Peak Demand} ($P$). Because physical server procurement requires long lead times, hardware sits idle during periods of average demand ($A$), generating waste.

Cloud computing transforms capital expenses (\textbf{CapEx}) into operational expenses (\textbf{OpEx}) through:
\begin{itemize}[leftmargin=*]
  \item \textbf{Common Infrastructure:} Standardized resource pools yielding statistical multiplexing gains.
  \item \textbf{Location Independence:} Aggregating demand across global time zones to smooth aggregate utilization curves.
  \item \textbf{Elasticity and Pay-As-You-Go:} Provisioning and releasing resources instantaneously to mirror real-time load.
\end{itemize}

\subsection{Formal Utility Premium Break-Even Analysis}
Let:
\begin{itemize}[leftmargin=*]
  \item $D(t)$: Instantaneous resource demand at time $t$ over time interval $[0, T]$.
  \item $P = \max_{0 \le t \le T} D(t)$: Peak resource demand.
  \item $A = \frac{1}{T} \int_{0}^{T} D(t)\, dt$: Average resource demand.
  \item $B$: Unit cost per unit time of owned (in-house baseline) infrastructure.
  \item $C$: Unit cost per unit time of utility cloud infrastructure.
  \item $U = \frac{C}{B}$: \textbf{Utility Premium} (the pricing multiple charged by the cloud provider over bare raw hardware ownership).
\end{itemize}

\begin{thmbox}{Cloud vs In-House Cost Break-Even Theorem}
Total in-house baseline cost ($B_T$) must provision for peak demand $P$:
\begin{equation}
B_T = P \times B \times T
\end{equation}
Total cloud cost ($C_T$) meters exact demand over time:
\begin{equation}
C_T = \int_{0}^{T} C \cdot D(t)\, dt = C \int_{0}^{T} D(t)\, dt = (U \cdot B) \cdot (A \cdot T) = A \cdot U \cdot B \cdot T
\end{equation}
Cloud computing is economically superior to owning physical hardware if and only if:
\begin{equation}
C_T < B_T \iff A \cdot U \cdot B \cdot T < P \cdot B \cdot T \iff U < \frac{P}{A}
\end{equation}
\textbf{Conclusion:} Cloud computing is cheaper than owning private hardware whenever the utility premium $U$ is strictly less than the \textbf{Peak-to-Average demand ratio} ($P/A$).
\end{thmbox}

\begin{figure}[H]
\centering
\begin{tikzpicture}
\begin{axis}[
    width=12cm, height=5.5cm,
    xlabel={Time $t$}, ylabel={Workload Demand $D(t)$},
    xmin=0, xmax=10, ymin=0, ymax=10,
    legend pos=north east,
    axis lines=left,
    xtick=\empty, ytick={3, 8},
    yticklabels={$A$ (Average), $P$ (Peak)}
]
\addplot[thick, navy, domain=0:10, samples=100] {5 + 3*sin(deg(x*1.5)) - cos(deg(x*0.5))};
\addlegendentry{Dynamic Demand $D(t)$}

\addplot[dashed, darkred, thick, domain=0:10] {8};
\addlegendentry{In-House Owned Capacity ($P$)}

\addplot[dotted, slate, thick, domain=0:10] {4.7};
\addlegendentry{Average Demand ($A$)}
\end{axis}
\end{tikzpicture}
\caption{Dynamic Demand vs In-House Static Over-Provisioning.}
\end{figure}

\subsection{Workload Penalty Integrals for Static Provisioning}
When an organization owns fixed resources $R(t)$, any divergence between available capacity $R(t)$ and instantaneous workload demand $D(t)$ incurs financial and operational penalties:
\begin{equation}
\text{Penalty Cost} \propto \int_{0}^{T} |D(t) - R(t)|\, dt
\end{equation}
\begin{itemize}[leftmargin=*]
  \item If $R(t) > D(t)$: The organization wastes capital paying for idle hardware, power, cooling, and software licensing.
  \item If $R(t) < D(t)$: The system suffers performance degradation, latency spikes, dropped transactions, and lost customer revenue.
\end{itemize}

Under exponential workload growth ($D(t) = e^t$), periodic static procurement with interval $t_p$ produces provisioning $R(t) = e^{t - t_p}$. The shortfall increases exponentially:
\begin{equation}
\text{Deficit}(t) = D(t) - R(t) = e^t - e^{t - t_p} = e^t (1 - e^{-t_p})
\end{equation}
demonstrating that fixed provisioning cycles can never keep pace with exponential demand growth.

\section{Big Data Architecture: Hadoop and HDFS}
\subsection{Hadoop Distributed File System (HDFS) Architecture}
HDFS is designed to store multi-terabyte and petabyte files reliably across large clusters of commodity hardware. It adopts a master/worker topology:
\begin{itemize}[leftmargin=*]
  \item \textbf{NameNode (Master):} Maintains the directory tree and file metadata in system RAM. File metadata maps filenames to a sequence of block IDs. Modifications are recorded durably in the \texttt{EditLog} on disk and consolidated periodically with the \texttt{FsImage} checkpoint file via the \textbf{Secondary NameNode}.
  \item \textbf{DataNodes (Workers):} Store actual data blocks as discrete files on local ext4/xfs filesystems. They handle block read/write operations from clients and emit heartbeat signals every 3 seconds to confirm node health.
\end{itemize}

\subsection{Block Size and Input Split Mathematics}
Unlike traditional filesystems with 4 KB disk clusters, HDFS employs very large block sizes (typically 64 MB, 128 MB, or 256 MB) to amortize disk seek times against sequential transfer rates.

When an input file of size $S$ is loaded into an HDFS cluster configured with block size $B$:
\begin{equation}
N_{\text{blocks}} = \left\lceil \frac{S}{B} \right\rceil
\end{equation}
In standard MapReduce configurations, Hadoop assigns exactly one \textbf{Mapper task} per logical \textbf{InputSplit} (which aligns 1:1 with HDFS physical blocks).

\begin{solved}{HDFS Mapper Task Allocation (Assignment 5, Q4)}
\Given An HDFS cluster operates with block size $B = 128\text{ MB}$. A client uploads an input file of size $S = 700\text{ MB}$ to execute a MapReduce job. Assuming default mapping of one map task per HDFS block, calculate the number of mapper tasks created.
\Formula
\begin{equation}
N_{\text{mappers}} = \left\lceil \frac{S}{B} \right\rceil
\end{equation}
\Steps
\begin{enumerate}[leftmargin=*]
  \item Decompose 700 MB across 128 MB block allocations:
  \begin{align*}
  \text{Block 1:} & \quad 128\text{ MB} & \text{Block 4:} & \quad 128\text{ MB}\\
  \text{Block 2:} & \quad 128\text{ MB} & \text{Block 5:} & \quad 128\text{ MB}\\
  \text{Block 3:} & \quad 128\text{ MB} & \text{Block 6:} & \quad 60\text{ MB} \quad (\text{Partial Block})
  \end{align*}
  \item Total allocated blocks: $5 \times 128\text{ MB} + 60\text{ MB} = 700\text{ MB} \implies 6\text{ blocks}$.
  \item Each block corresponds to an input split: $N_{\text{mappers}} = 6$.
\end{enumerate}
\Answer Exactly \textbf{6 mapper tasks} are spawned.
\end{solved}

\subsection{MapReduce Execution Pipeline: Partitioner vs Combiner}
The lecture highlights two critical intermediate optimization components in MapReduce:
\begin{itemize}[leftmargin=*]
  \item \textbf{Combiner (Mini-Reducer):} An \emph{optional} optimization component executed on the map worker node prior to network shuffle. It performs local pre-aggregation on intermediate key-value pairs (e.g., computing sub-totals in a word count job), drastically decreasing the volume of intermediate data transmitted across the cluster network. The execution of a Combiner is never guaranteed; Hadoop may run it zero, one, or multiple times depending on memory pressure.
  \item \textbf{Partitioner:} A \emph{mandatory} component that determines which physical Reducer instance receives each intermediate key. By default, Hadoop utilizes the \textbf{HashPartitioner}:
  \begin{equation}
  \text{Reducer ID} = (\text{key.hashCode()} \ \& \ \text{Integer.MAX\_VALUE}) \pmod{N_{\text{reducers}}}
  \end{equation}
  This guarantees that all intermediate values associated with the identical key arrive at the exact same Reducer node.
\end{itemize}

\section{Green Cloud Computing and Server Energy Dynamics}
\subsection{Data Center Energy Footprint and Efficiency Metrics}
Massive data centers consume hundreds of megawatts of electricity. Energy consumption is divided into two broad categories:
\begin{enumerate}[leftmargin=*]
  \item \textbf{IT Equipment Power ($P_{\text{IT}}$):} Compute servers, storage arrays, and network switches.
  \item \textbf{Infrastructure Overhead Power ($P_{\text{overhead}}$):} Chiller plants, cooling towers, fans, uninterruptible power supplies (UPS), power distribution units (PDUs), and facility lighting.
\end{enumerate}

\begin{defbox}{Power Usage Effectiveness (PUE) and DCiE}
The industry standard metric established by The Green Grid is \textbf{Power Usage Effectiveness (PUE)}:
\begin{equation}
\text{PUE} = \frac{\text{Total Facility Power}}{\text{IT Equipment Power}} = \frac{P_{\text{IT}} + P_{\text{overhead}}}{P_{\text{IT}}} \ge 1.0
\end{equation}
A PUE of $2.0$ indicates that for every 1 Watt powering servers, another 1 Watt is consumed by cooling and power delivery. Modern hyperscale facilities (Google, Microsoft) achieve PUE values between $1.1$ and $1.2$.\\
The inverse metric is \textbf{Data Center Infrastructure Efficiency (DCiE)}:
\begin{equation}
\text{DCiE} = \frac{1}{\text{PUE}} \times 100\%
\end{equation}
\end{defbox}

\subsection{Server Power Dynamics and Core Utilization}
A physical server's power consumption is not static; it scales with CPU utilization and active hardware components:
\begin{equation}
P(u) = P_{\text{idle}} + (P_{\text{peak}} - P_{\text{idle}}) \times u
\end{equation}
where $P_{\text{idle}}$ is the idle power draw (typically $50\%$ to $60\%$ of peak power $P_{\text{peak}}$), and $u \in [0, 1]$ represents CPU utilization.

\begin{thmbox}{Impact of Dynamic Core Activation on Power Draw}
When a multi-core server detects a workload spike and dynamically switches CPU cores from low-power C-states (sleep states) to active C0 execution states:
\begin{enumerate}[leftmargin=*]
  \item Each active core draws dynamic switching power: $P_{\text{dynamic}} = \alpha \cdot C_L \cdot V^2 \cdot f$ (where $\alpha$ is activity factor, $C_L$ load capacitance, $V$ operating voltage, and $f$ clock frequency).
  \item Static leakage current power increases proportionally with core temperature.
  \item \textbf{Net Effect:} Total server power consumption \textbf{increases} as additional CPU cores become active, even if workload execution time per core drops.
\end{enumerate}
\end{thmbox}

\section{Comprehensive Solved Assignment Problems (Week 5)}

\begin{solved}{SLA Downtime Allowance Calculation (Assignment 5, Q1)}
\Given A cloud provider commits to a 99\% uptime SLA. The client's application runs 12 hours/day for 30 days. Find maximum allowable downtime.
\Steps
\begin{equation}
T_{\text{scheduled}} = 12\text{ h/day} \times 30\text{ days} = 360\text{ hours}
\end{equation}
\begin{equation}
T_{\text{allowed downtime}} = (1 - 0.99) \times 360 = 0.01 \times 360 = 3.6\text{ hours}
\end{equation}
\Answer \textbf{Option A: 3.6 hours}.
\end{solved}

\begin{solved}{Multi-Core Server Power Consumption (Assignment 5, Q2)}
\Given A cloud server dynamically activates additional CPU cores to handle a sudden surge in user requests.
\Steps
\begin{enumerate}[leftmargin=*]
  \item Activating dormant processor cores draws additional dynamic switching power and leakage current.
  \item While techniques like DVFS reduce voltage at lower frequencies, adding active cores increases total aggregate wattage drawn by the system.
\end{enumerate}
\Answer \textbf{Option C: Power consumption increases as additional CPU cores become active}.
\end{solved}

\begin{solved}{Mitigating Upfront Capital Expenditure (Assignment 5, Q3)}
\Given Which cloud characteristic directly eliminates massive upfront hardware purchases?
\Steps
\begin{enumerate}[leftmargin=*]
  \item \textbf{Pay-as-you-go pricing} converts Capital Expenditure (CapEx) into Operational Expenditure (OpEx), allowing organizations to pay strictly for consumed compute time and storage without purchasing servers upfront.
\end{enumerate}
\Answer \textbf{Option B: Pay-as-you-go pricing}.
\end{solved}

\begin{solved}{HDFS Block Input Splits and Map Tasks (Assignment 5, Q4)}
\Given Block size $B = 128\text{ MB}$, input file $S = 700\text{ MB}$.
\Steps
\begin{equation}
N = \left\lceil \frac{700}{128} \right\rceil = \left\lceil 5.46875 \right\rceil = 6
\end{equation}
\Answer \textbf{Option B: 6}.
\end{solved}

\begin{solved}{Primary Virtualization Management Entity (Assignment 5, Q5)}
\Given Identify the virtualization layer component responsible for creating, running, and managing virtual machines.
\Steps
\begin{enumerate}[leftmargin=*]
  \item The \textbf{Hypervisor} (or Virtual Machine Monitor, VMM) abstracts physical hardware resources and mediates guest virtual machine execution.
\end{enumerate}
\Answer \textbf{Option B: Hypervisor}.
\end{solved}

\begin{solved}{MapReduce Key-to-Reducer Assignment (Assignment 5, Q6)}
\Given Identify the MapReduce subsystem that routes each intermediate key to its designated Reducer.
\Steps
\begin{enumerate}[leftmargin=*]
  \item The \textbf{Partitioner} calculates the destination reducer index for each key (typically via hash modulo operations).
\end{enumerate}
\Answer \textbf{Option B: Partitioner}.
\end{solved}

\begin{solved}{Optional MapReduce Optimization Phase (Assignment 5, Q7)}
\Given Identify the MapReduce component that is optional and may execute zero, one, or multiple times.
\Steps
\begin{enumerate}[leftmargin=*]
  \item The \textbf{Combiner} operates as a semi-reducer on local map output. Hadoop does not guarantee its execution; therefore, MapReduce job logic must remain algebraically valid even when the Combiner is skipped entirely.
\end{enumerate}
\Answer \textbf{Option C: Combiner}.
\end{solved}

\begin{solved}{Data Center Energy Consumption Contributors (Assignment 5, Q8)}
\Given Which data center subsystems contribute directly to overall energy consumption?
\Steps
\begin{enumerate}[leftmargin=*]
  \item IT server equipment (processors, memory, disks).
  \item Air conditioning chillers, cooling towers, and ventilation fans.
  \item Power distribution units (PDUs) and uninterruptible power supply (UPS) transformers.
\end{enumerate}
\Answer \textbf{Option D: All of the above}.
\end{solved}

\begin{solved}{Server Consolidation Impact on Green Computing (Assignment 5, Q9)}
\Given How does server consolidation via live VM migration promote green cloud computing?
\Steps
\begin{enumerate}[leftmargin=*]
  \item Packing virtual machines onto fewer physical hosts allows underutilized hosts to be powered down into deep sleep states, eliminating idle power waste.
\end{enumerate}
\Answer \textbf{Option A: Reduces total energy consumption by powering down idle physical servers}.
\end{solved}

\begin{solved}{HDFS Data Integrity and Fault Tolerance (Assignment 5, Q10)}
\Given How does HDFS achieve fault tolerance against DataNode drive crashes without relying on RAID hardware?
\Steps
\begin{enumerate}[leftmargin=*]
  \item HDFS replicates each data block across multiple physically independent DataNodes (default replication factor $R = 3$) with rack awareness.
\end{enumerate}
\Answer \textbf{Option B: Block Replication across multiple nodes}.
\end{solved}
